\scanpage{12}
sont des groupes fondamentaux ``géométriques'' --- \unsure{qui|que} peuvent d'ailleurs se calculer par voie transcendante, en utilisant un plongement de $\overline{\mathbb Q}$ dans $\mathbb C$ (d'où un isom $\overline{\mathbb Q}\simeq\overline{\mathbb Q}_0$), comme les complétés profinis des $\pi_1(\unsure{U_i(\underset{k}{\mathbb C})|U_i(\mathbb C)},\bar\eta)$, \unsure{où|en} maintenant $\bar\eta$ est interprété comme un pt \deleted{complexe} commun des espaces top. sous-jacents aux variétés complexes $U_i(\mathbb C)$.

La suite exacte (8) est donc la limite projective des suites exactes d'homotopie (12), ce qui donne en particulier
\[
\begin{aligned}
(13)\qquad \pi_{\bar K/K}&\simeq\varprojlim_i\underbrace{\pi_1(\bar U_i,\bar\eta)}_{\pi_i}\\
&\simeq\varprojlim_i\widehat{\pi_1(U_i(\mathbb C),\bar\eta)}
\end{aligned}
\]
\margin{Cette interprétation n'est évidemment pas particulière au corps de base $\mathbb Q$ --- elle vaut \unsure{pour|pr} tte ext. de corps $K/K_0$}

Utilisant les fibrations des $U_i$ (dans le cas où on s'astreint à prendre des variétés élémentaires d'Artin), on trouve que tout $\pi_i$ est un groupe \deleted{\unsure{par}} extension successive de \inserted{$n$} groupes profinis \uline{libres} (ici $n=\dim U_i=\operatorname{deg.tr.}K/\mathbb Q$). Ceci redonne
\margin{11}
